% ============================================================
% ABSTRACT
% ============================================================
\begin{abstract}
Test-time compute improves decision quality for LLM agents, but deciding when it is beneficial remains unresolved. Existing gating methods rely on uncertainty signals, implicitly assuming a fixed relationship between signal value and compute benefit. We show that this assumption is incorrect. 
Across eight diverse agent environments and three model backbones, we find that the same signal value that predicts rollout benefit in one environment predicts rollout harm in another, indicating a reversal in the signal--utility relationship. 
We trace this to a conflation of two uncertainty sources, \emph{information poverty} and \emph{decision difficulty}, and formalize it via a two-source model grounded in Simpson's paradox, where subgroup-level trends reverse upon aggregation. We prove that direction discovery is a necessary condition for non-negative value of computation across environments. 
Motivated by this result, we propose EAAG (Environment-Aware Adaptive Gating), which learns environment-specific signal--utility relationships through direction discovery. Across eight environments, EAAG outperforms all fixed-direction baselines with zero additional per-step overhead, achieving adaptive gating behavior solely by learning signal direction. 
\end{abstract}
